Nine baselines (Fig.~\ref{fig:stage1}) anchor the design space. Pristine zigzag graphene (40\,\AA{} cell, homogeneous) reaches $\sigma_\mathrm{max}=38.8$\,N/m at 27\% strain and fails in a single avalanche; the dip in its curve near 18\% is the transverse-branch switch of the bistable homogeneous response discussed in Sec.~\ref{sec:validation}. Pristine armchair graphene peaks earlier and lower ($30.7$\,N/m at 18\%) and then softens continuously without bond loss until rupture at 30\% (the model's armchair response). The 2D modulus of all dense sheets is 224--250\,N/m ($287$\,N/m for armchair in the uniaxial-stress protocol), consistent with the elastic validation. Two per cent vacancies reduce the strength to $\approx24$\,N/m whether distributed (23.6\,N/m) or clustered (24.4\,N/m), but the fracture mode differs: distributed vacancies fail progressively (first bond loss at 9.8\%, peak at 12.8\%, failure at 14.3\%, six damage steps, load retained after the peak), while the clustered ones fail in one avalanche after a single first damage event. Central slit cracks of 10, 20 and 30\,\AA{} give 24.1, 18.5 and 18.7\,N/m; the 10$\to$20\,\AA{} step follows a Griffith-type power law (fitted exponent $-0.38$ against $-0.5$), but the 30\,\AA{} crack is \emph{not} weaker than the 20\,\AA{} crack: it starts to lose bonds at 7.9\% strain, grows stably bond by bond while the load still rises to 9.9\%, and only then runs. A circular hole of 20\,\AA{} (19.1\,N/m) is as damaging as a 20\,\AA{} crack. The reference nanomesh (period 16\,\AA, $\phi=0.22$) has $Y_{2D}=119$\,N/m, $\sigma_\mathrm{max}=12.0$\,N/m at 14\%, first damage at 10\% and fails progressively over 17\% strain with $W=1.1$\,J/m$^2$, i.e. it is 3$\times$ weaker than pristine graphene but retains load over a strain range comparable to the intact sheet's elastic range. Two features that recur later are already visible: (i) architecture changes the \emph{mode} of fracture at nearly constant strength (vacancy organisation), and (ii) sharp flaws beyond $\approx$20\,\AA{} enter a flaw-tolerant regime in this model, where stable crack growth precedes the instability.
